% GENERATED FILE. Edit canonical lessons, then run npm run generate:books.
% canonical-source-hash: fnv1a64:1c7df633ec8de036
% canonical-lessons: TA-C460-viyappu, TA-C460-malaippu, TA-C460-tikai, TA-C460-paravacam, TA-C460-cilir, TA-R460-review-453-460

\chapter{Wonder and Awe}
\label{ch:ta-wonder-and-awe}
\hlchaptermodality{460}

\begin{chapteropening}
\textbf{By the end of this chapter:} \emph{I can talk about amazement, awe and being thrilled.}

Complete the last lesson of chapter 460: I can talk about amazement, awe and being thrilled.
\end{chapteropening}

\section[\texorpdfstring{\ta{வியப்பு} (viyappu)}{வியப்பு (viyappu)}]{\ta{வியப்பு} (viyappu) — amazement; wonder}

\label{lesson:TA-C460-viyappu}



\begin{quote}
Before the new one: say the Tamil for glee, then the Tamil for merriment.
\end{quote}

\subsection*{You'll want to know: \ta{வியப்பு}}
\textbf{\ta{வியப்பு}} — \emph{viyappu} — \textquotedblleft{}amazement; wonder\textquotedblright{}.

\begin{quote}
\ta{குழந்தையின்} \ta{திறமை} \ta{எனக்கு} \ta{வியப்பாக} \ta{இருந்தது}.
\end{quote}

\begin{quote}
\emph{kuḻantaiyiṉ tiṟamai eṉakku viyappāka iruntatu.}
\end{quote}

\begin{quote}
The child’s talent amazed me.
\end{quote}

The verb is \ta{விய} \textquotedblleft{}to wonder\textquotedblright{}, with its past in -\ta{ந்த}-: \ta{வியந்தேன்}, I marvelled.

\begin{cousinweb}
\textbf{\ta{வியப்பு}} (\emph{viyappu}) is native Dravidian, a noun from \ta{விய} \emph{viya}, \textquotedblleft{}to wonder, to admire\textquotedblright{} (DEDR 5403).
\end{cousinweb}

\begin{grammarlens}[title={What you've built}]
One more word for describing what you lived and felt.
\end{grammarlens}

\subsection*{Your turn}
\begin{itemize}
\raggedright
  \item \emph{Say it:} \emph{viyappu}
  \item \emph{Say it:} \emph{viyappu}, once more
  \item \emph{Say it:} \emph{kuḻantaiyiṉ tiṟamai eṉakku viyappāka iruntatu.} — The child’s talent amazed me.
  \item \emph{Recall:} say the Tamil for glee, then the Tamil for merriment, then say \emph{viyappu} again
\end{itemize}

\begin{tcolorbox}[breakable,colback=teal!4,colframe=teal!35!black,title={Before you move on}]
What does \ta{வியப்பு} mean? (\textquotedblleft{}Amazement; wonder\textquotedblright{}.) Say it once more.
\end{tcolorbox}
\section[\texorpdfstring{\ta{மலைப்பு} (malaippu)}{மலைப்பு (malaippu)}]{\ta{மலைப்பு} (malaippu) — awe; being overwhelmed}

\label{lesson:TA-C460-malaippu}



\begin{quote}
Before the new one: say the Tamil for merriment, then the Tamil for amazement.
\end{quote}

\subsection*{You'll want to know: \ta{மலைப்பு}}
\textbf{\ta{மலைப்பு}} — \emph{malaippu} — \textquotedblleft{}awe; being overwhelmed\textquotedblright{}.

\begin{quote}
\ta{அவ்வளவு} \ta{பெரிய} \ta{கோயிலைப்} \ta{பார்த்து} \ta{எனக்கு} \ta{மலைப்பாக} \ta{இருந்தது}.
\end{quote}

\begin{quote}
\emph{avvaḷavu periya kōyilaip pārttu eṉakku malaippāka iruntatu.}
\end{quote}

\begin{quote}
I was awestruck to see such a big temple.
\end{quote}

Not from \ta{மலை} \textquotedblleft{}a hill\textquotedblright{} but from a verb \ta{மலை} \textquotedblleft{}to be staggered\textquotedblright{}.

\begin{cousinweb}
\textbf{\ta{மலைப்பு}} (\emph{malaippu}) is native Dravidian (DEDR 4736): the entry gives a verb \ta{மலை} \emph{malai}, \textquotedblleft{}to be staggered, amazed\textquotedblright{}, and this noun beside it.
\end{cousinweb}

\begin{grammarlens}[title={What you've built}]
One more word for describing what you lived and felt.
\end{grammarlens}

\subsection*{Your turn}
\begin{itemize}
\raggedright
  \item \emph{Say it:} \emph{malaippu}
  \item \emph{Say it:} \emph{malaippu}, once more
  \item \emph{Say it:} \emph{avvaḷavu periya kōyilaip pārttu eṉakku malaippāka iruntatu.} — I was awestruck to see such a big temple.
  \item \emph{Recall:} say the Tamil for merriment, then the Tamil for amazement, then say \emph{malaippu} again
\end{itemize}

\begin{tcolorbox}[breakable,colback=teal!4,colframe=teal!35!black,title={Before you move on}]
What does \ta{மலைப்பு} mean? (\textquotedblleft{}Awe; being overwhelmed\textquotedblright{}.) Say it once more.
\end{tcolorbox}
\section[\texorpdfstring{\ta{திகை} (tikai)}{திகை (tikai)}]{\ta{திகை} (tikai) — to be stunned; to be bewildered}

\label{lesson:TA-C460-tikai}



\begin{quote}
Before the new one: say the Tamil for amazement, then the Tamil for awe.
\end{quote}

\subsection*{You'll want to know: \ta{திகை}}
\textbf{\ta{திகை}} — \emph{tikai} — \textquotedblleft{}to be stunned; to be bewildered\textquotedblright{}.

\begin{quote}
\ta{அவனுடைய} \ta{பதிலைக்} \ta{கேட்டு} \ta{நான்} \ta{திகைத்தேன்}.
\end{quote}

\begin{quote}
\emph{avaṉuṭaiya patilaik kēṭṭu nāṉ tikaittēṉ.}
\end{quote}

\begin{quote}
I was stunned to hear his answer.
\end{quote}

Past (I …): \ta{திகைத்தேன்}, \emph{tikaittēṉ}.

\ta{திகைத்து} \ta{நின்றேன்}, I stood dumbfounded.

\begin{cousinweb}
\textbf{\ta{திகை}} (\emph{tikai}) is native Dravidian (DEDR 3207); the entry also gives \ta{திகைப்பு} \emph{tikaippu}, \textquotedblleft{}amazement\textquotedblright{}.
\end{cousinweb}

\begin{grammarlens}[title={What you've built}]
One more word for describing what you lived and felt.
\end{grammarlens}

\subsection*{Your turn}
\begin{itemize}
\raggedright
  \item \emph{Say it:} \emph{tikai}
  \item \emph{Say it:} \emph{tikai}, once more
  \item \emph{Say it:} \emph{avaṉuṭaiya patilaik kēṭṭu nāṉ tikaittēṉ.} — I was stunned to hear his answer.
  \item \emph{Recall:} say the Tamil for amazement, then the Tamil for awe, then say \emph{tikai} again
\end{itemize}

\begin{tcolorbox}[breakable,colback=teal!4,colframe=teal!35!black,title={Before you move on}]
What does \ta{திகை} mean? (\textquotedblleft{}To be stunned; to be bewildered\textquotedblright{}.) Say it once more.
\end{tcolorbox}
\section[\texorpdfstring{\ta{பரவசம்} (paravacam)}{பரவசம் (paravacam)}]{\ta{பரவசம்} (paravacam) — ecstasy; rapture}

\label{lesson:TA-C460-paravacam}



\begin{quote}
Before the new one: say the Tamil for awe, then the Tamil for to be stunned.
\end{quote}

\subsection*{You'll want to know: \ta{பரவசம்}}
\textbf{\ta{பரவசம்}} — \emph{paravacam} — \textquotedblleft{}ecstasy; rapture\textquotedblright{}.

\begin{quote}
\ta{இசையைக்} \ta{கேட்டு} \ta{அவர்} \ta{பரவசம்} \ta{அடைந்தார்}.
\end{quote}

\begin{quote}
\emph{icaiyaik kēṭṭu avar paravacam aṭaintār.}
\end{quote}

\begin{quote}
Listening to the music, he was enraptured.
\end{quote}

\ta{பரவசம்} \ta{அடை}, to be carried away.

\begin{cousinweb}
\textbf{\ta{பரவசம்}} (\emph{paravacam}) is borrowed from Sanskrit \emph{paravaśa}, \textquotedblleft{}under another's power\textquotedblright{} (Monier-Williams): in Tamil, carried away by feeling.
\end{cousinweb}

\begin{grammarlens}[title={What you've built}]
One more word for describing what you lived and felt.
\end{grammarlens}

\subsection*{Your turn}
\begin{itemize}
\raggedright
  \item \emph{Say it:} \emph{paravacam}
  \item \emph{Say it:} \emph{paravacam}, once more
  \item \emph{Say it:} \emph{icaiyaik kēṭṭu avar paravacam aṭaintār.} — Listening to the music, he was enraptured.
  \item \emph{Recall:} say the Tamil for awe, then the Tamil for to be stunned, then say \emph{paravacam} again
\end{itemize}

\begin{tcolorbox}[breakable,colback=teal!4,colframe=teal!35!black,title={Before you move on}]
What does \ta{பரவசம்} mean? (\textquotedblleft{}Ecstasy; rapture\textquotedblright{}.) Say it once more.
\end{tcolorbox}
\section[\texorpdfstring{\ta{சிலிர்} (cilir)}{சிலிர் (cilir)}]{\ta{சிலிர்} (cilir) — to tingle; to get goosebumps}

\label{lesson:TA-C460-cilir}



\begin{quote}
Before the new one: say the Tamil for to be stunned, then the Tamil for ecstasy.
\end{quote}

\subsection*{You'll want to know: \ta{சிலிர்}}
\textbf{\ta{சிலிர்}} — \emph{cilir} — \textquotedblleft{}to tingle; to get goosebumps\textquotedblright{}.

\begin{quote}
\ta{அந்தப்} \ta{பாடலைக்} \ta{கேட்டதும்} \ta{என்} \ta{உடல்} \ta{சிலிர்த்தது}.
\end{quote}

\begin{quote}
\emph{antap pāṭalaik kēṭṭatum eṉ uṭal cilirttatu.}
\end{quote}

\begin{quote}
When I heard that song, my body tingled.
\end{quote}

Past (it …): \ta{சிலிர்த்தது}, \emph{cilirttatu}.

The body or the hair is the subject, so the past is \ta{சிலிர்த்தது}.

\begin{cousinweb}
\textbf{\ta{சிலிர்}} (\emph{cilir}) is native Dravidian (DEDR 2565): hair standing on end from fright or delight.
\end{cousinweb}

\begin{grammarlens}[title={What you've built}]
One more word for describing what you lived and felt.
\end{grammarlens}

\subsection*{Your turn}
\begin{itemize}
\raggedright
  \item \emph{Say it:} \emph{cilir}
  \item \emph{Say it:} \emph{cilir}, once more
  \item \emph{Say it:} \emph{antap pāṭalaik kēṭṭatum eṉ uṭal cilirttatu.} — When I heard that song, my body tingled.
  \item \emph{Recall:} say the Tamil for to be stunned, then the Tamil for ecstasy, then say \emph{cilir} again
\end{itemize}

\begin{tcolorbox}[breakable,colback=teal!4,colframe=teal!35!black,title={Before you move on}]
What does \ta{சிலிர்} mean? (\textquotedblleft{}To tingle; to get goosebumps\textquotedblright{}.) Say it once more.
\end{tcolorbox}
\section[(dialogue)]{Looking back at feelings}

\label{lesson:TA-R460-review-453-460}



\begin{quote}
Say the Tamil for rage and irritation.
\end{quote}

\subsection*{Your turn}
Say these aloud:

\begin{itemize}
\raggedright
  \item rage, irritation, to flare up, to seethe, to mutter
  \item terror, panic, shock, to be startled, to tremble
  \item unease, anguish, frustration, to bemoan, to suffer
  \item breath, a sigh, relief, consolation, to subside
  \item mercy, tenderness, sympathy, pitifulness, concern
  \item to wail, to sob, to whimper, to lament, to scream
  \item bliss, pleasure, elation, glee, merriment
  \item amazement, awe, to be stunned, ecstasy, to tingle
\end{itemize}

\begin{tcolorbox}[breakable,colback=teal!4,colframe=teal!35!black,title={Before you move on}]
Rage? (\textbf{\ta{ஆத்திரம்}}, \emph{āttiram}.) Irritation? (\textbf{\ta{எரிச்சல்}}, \emph{ericcal}.) To flare up? (\textbf{\ta{சீறு}}, \emph{cīṟu}.) To seethe? (\textbf{\ta{கொந்தளி}}, \emph{kontaḷi}.) To mutter? (\textbf{\ta{முணுமுணு}}, \emph{muṇumuṇu}.) Terror? (\textbf{\ta{திகில்}}, \emph{tikil}.) Panic? (\textbf{\ta{பீதி}}, \emph{pīti}.) Shock? (\textbf{\ta{அதிர்ச்சி}}, \emph{atircci}.) To be startled? (\textbf{\ta{திடுக்கிடு}}, \emph{tiṭukkiṭu}.) To tremble? (\textbf{\ta{நடுங்கு}}, \emph{naṭuṅku}.) Unease? (\textbf{\ta{சஞ்சலம்}}, \emph{cañcalam}.) Anguish? (\textbf{\ta{உளைச்சல்}}, \emph{uḷaiccal}.) Frustration? (\textbf{\ta{ஆதங்கம்}}, \emph{ātaṅkam}.) To bemoan? (\textbf{\ta{அங்கலாய்}}, \emph{aṅkalāy}.) To suffer? (\textbf{\ta{தவி}}, \emph{tavi}.) Breath? (\textbf{\ta{மூச்சு}}, \emph{mūccu}.) A sigh? (\textbf{\ta{பெருமூச்சு}}, \emph{perumūccu}.) Relief? (\textbf{\ta{நிம்மதி}}, \emph{nimmati}.) Consolation? (\textbf{\ta{ஆறுதல்}}, \emph{āṟutal}.) To subside? (\textbf{\ta{தணி}}, \emph{taṇi}.) Mercy? (\textbf{\ta{கருணை}}, \emph{karuṇai}.) Tenderness? (\textbf{\ta{பரிவு}}, \emph{parivu}.) Sympathy? (\textbf{\ta{அனுதாபம்}}, \emph{aṉutāpam}.) Pitifulness? (\textbf{\ta{பரிதாபம்}}, \emph{paritāpam}.) Concern? (\textbf{\ta{அக்கறை}}, \emph{akkaṟai}.) To wail? (\textbf{\ta{கதறு}}, \emph{kataṟu}.) To sob? (\textbf{\ta{விம்மு}}, \emph{vimmu}.) To whimper? (\textbf{\ta{விசும்பு}}, \emph{vicumpu}.) To lament? (\textbf{\ta{புலம்பு}}, \emph{pulampu}.) To scream? (\textbf{\ta{அலறு}}, \emph{alaṟu}.) Bliss? (\textbf{\ta{ஆனந்தம்}}, \emph{āṉantam}.) Pleasure? (\textbf{\ta{இன்பம்}}, \emph{iṉpam}.) Elation? (\textbf{\ta{பூரிப்பு}}, \emph{pūrippu}.) Glee? (\textbf{\ta{குதூகலம்}}, \emph{kutūkalam}.) Merriment? (\textbf{\ta{உல்லாசம்}}, \emph{ullācam}.) Amazement? (\textbf{\ta{வியப்பு}}, \emph{viyappu}.) Awe? (\textbf{\ta{மலைப்பு}}, \emph{malaippu}.) To be stunned? (\textbf{\ta{திகை}}, \emph{tikai}.) Ecstasy? (\textbf{\ta{பரவசம்}}, \emph{paravacam}.) To tingle? (\textbf{\ta{சிலிர்}}, \emph{cilir}.)
\end{tcolorbox}
